\documentclass{article}
\usepackage{graphicx} % Required for inserting images
\usepackage{amsmath}
\usepackage{amssymb}
\title{Polynomial}
\author{Sun Shengran}
\date{December 2025}

\begin{document}

\maketitle

\section{Introduction to Polynomial}
\subsection{Field}

Let F be the set of some numbers.
If F is closed to addition, subtraction, multiplication and division, then F is called a field.

\noindent
For example: rational set $\mathbb{Q}$ (Quotient) $\rightarrow$\ rational field

\subsection{Polynomial}
\subsubsection{Definition}
For a field F and a notation(variation) \textit{X} (probably not $\in$ F),

\begin{center}
$f(X) = a_n X^n + a_{n-1} X^{n-1} + ... + a_1 X + a_0$ where $a_i \in F,\ i = 1,2,...,n$
\end{center}

\noindent
is a polynomial with single variation on F.
\subsubsection{Basic Features}
1. If $n \ge 0$ and $a_n \neq 0$, then $n = \deg f(x)$.

\vspace{0.2cm}
\noindent
2. Any $a \neq 0 \in F$ can be seen as a polynomial with degree 0.

\vspace{0.2cm}
\noindent
3. $0 \in F$ is a polynomial with degree $-\infty$ and with no first term.

\vspace{0.2cm}
\noindent
4. For polynomials
$f(X) = a_nX^n +...+a_0$ and $g(X) = b_mX^m +...+ b_0$,

\begin{center}
$f(X) = g(x)$ $\Leftrightarrow$\ $m = n$ and $a_i = b_i,\ i = 1,2,...,n$
\end{center}

\vspace{0.2cm}
\noindent
5. $\deg fg = \deg f + \deg g$

\vspace{0.2cm}
\noindent
6. $f \neq 0,\ g \neq 0$ $\Leftrightarrow$\ $fg \neq 0$ $\Leftrightarrow$\ $F[X]$ has no null factor.

\vspace{0.2cm}
\noindent
7. Most polynomials $f(X)$ on F are non-invertible. The only invertible polynomial is non-zero constants(F*).

\section{Number Theory in Polynomial}
\subsection{Euclidean Division}
$\forall f, g \in F[X]$ and $g \neq 0$,\  $\exists ! \ q, r \in F[X]$ such that

\begin{center}
$f = gq + r\ (r = 0\  \text{or}\  \deg r < \deg g)$
\end{center}

\subsection{Factors of $f(X)$}
\subsubsection{Factor}
If $f = gq\ (g \neq 0)$, then $g|f$ so $g$ is a factor of $f$.\ 

\vspace{0.2cm}
\noindent
Trivial factors: 1 and $f(X)$ itself.

\subsubsection{Reducibility}
\noindent
Also, if $f = gq$ and $g,q\in F$ are not constants, then $f$ is reducible over F.

\vspace{0.2cm}
\noindent
If $f(X)$ is irreducible (i.e. only has trivial factors), then $f$ is prime.

\subsubsection{Association}
\noindent
If $g|f$, then $bg|af$ for any non-zero constant $a,b \in F$.

\vspace{0.2cm}
\noindent
For any non-zero constant $c \in F$, $g$ and $cg$ are associate.

\vspace{0.2cm}
\noindent
If $g|f$ and $f|g$, then $f$ and $g$ are associate.

\subsubsection{Highest Common Factor}
\noindent
\begin{center}
$d(X)$ is the HCF of $f(X)$ and $g(X)$

$\Updownarrow$\

$d|f, d|g$\  and\  $\forall c \in F,\  c|f,\  c|g,\  c|d$
\end{center}

\noindent
The HCF of $f$ and $g$ is unique if associate is not considered.\\
If $(f,g) = 1$, then $f$ and $g$ are coprime.

\subsection{B\'{e}zout Identity}
For $f(X)$ and $g(X) \in F[X]$, their HCF $d(X)$ always exists uniquely\ 
and there exists $u(X), v(X) \in F[X]$\ 
such that

\begin{center}
$u(X)f(X) + v(X)g(X) =  d(X)$
\end{center}

\noindent
Particularly,

\noindent
\begin{center}
$f(X)$ and $g(X)$ are coprime

$\Updownarrow$

$\exists\  u(X), v(X) \in F[X]$ such that $u(X)f(X) + v(X)g(X) = 1$
\end{center}

\noindent
Generally,

\noindent
For $f_i \in F[X],\ i = 1,2,...,s$, $\exists\ u_i \in F[X],\ i=1,2,...s$ such that
\begin{center}
$u_1f_1+u_2f_2+...+u_sf_s = (f_1,f_2,...f_s)$
\end{center}

\subsection{Unique Factorisation Domain}
$\forall$ non-constant $f(X) \in F[X]$ can be factorised uniquely:

\begin{center}
$f(X) = p_1(X)^{v_1} p_2(X)^{v_2}... p_r(X)^{v_r}$
where $p_i(X)$ are prime polynomials in $F[X]$.

\vspace{0.15cm}
$p_i(X)^{v_i}||f(X)$
\end{center}

\noindent
This theorem can be stated as $F[X]$ is a unique factorisation domain.

\subsection{Remainder Theorem and Factor Theorem}
Remainder Theorem: $f(X)$ has remainder $r=f(c)$ when divided by $(X-c)$

\noindent
Factor Theorem: $(X-c)|f(X)$\  $\Leftrightarrow$\ $f(c)=0$

\vspace{0.1cm}
\noindent
If $f(X)=(X-c)^m g(X)\ \ (g(c)\neq 0, m \ge 1)$, c is called multiple root when $m \ge 2$,\ 
c is called single root when $m=1$.

\subsection{Foundamental Theorem of Algebra}
\textbf{Theorem}: If $f(X)$ is of degree $n (n \ge 0)$ in $F[X]$, then $f(X)$ has at most $n$ roots over $F$.

\noindent
(Proof by induction)

\noindent
Attention: If we deal with a set which is not a field,\ 
then a polynomial with degree $n$ could have $>n$ roots. 

\vspace{0.2cm}
\noindent
\textbf{Corollary}: If $f, g \in F[X]$ are both of degree $\le n (n \ge 0)$ and have $n+1$ points with equal values, then $f=g$.

\noindent
(Proof: consider $h = f - g$)

\subsection{Differentiation}
To justify the multiple roots of polynomial, we introduce the concept of differentiation.

\subsubsection{Definition}
The differentiation form of $f(X)= a_n X^n + a_{n-1} X^{n-1} +...+ a_1 X + a_0$ is

\noindent
\begin{center}
$f'(X) = na_n X^{n-1} + (n-1)a_{n-1} X^{n-2} +...+ a_1 \in F[X]$
\end{center}

\subsubsection{Theorem}
For $f(X)\in F[X]$, $c \in F$,

\vspace{0.1cm}
\noindent
\textbf{Theorem}: $c$ is the root of $f(X)$ $\Leftrightarrow$\ $f(c) = f'(c) = 0$
$\Leftrightarrow$\ $c$ is the root of $(f, f')$

\vspace{0.1cm}
\noindent
\textbf{Corollary1}: If $(f,f')=1$, then $f(X)$ has no multiple root in any field $E \supseteq F$

\vspace{0.1cm}
\noindent
\textbf{Corollary2}: If $f(X)$ is irreducible in $F \subset \mathbb{C}$, then $f(X)$ has no multiple root.

\subsection{Summary}
The application of number theory in polynomial is part of algebraic number theory.
\end{document}
